\chapter{Атомы, связи и твёрдое тело}\label{ch:atoms}

Почти все свойства металла --- температура плавления, жёсткость, тепловое расширение, растворимость водорода ---
выводятся из одного: как энергия пары соседних атомов зависит от расстояния между ними и что при этом делают
электроны. Эта глава даёт минимум физики твёрдого тела, нужный, чтобы последующие главы не казались набором
эмпирических формул.

\section{Атом и его внешние электроны}

Атом --- ядро и электроны на оболочках. Химию и связь определяют \textbf{внешние (валентные)} электроны; внутренние
оболочки заполнены и в связи почти не участвуют. Порядок заполнения оболочек задаётся правилом
$1s\,2s\,2p\,3s\,3p\,4s\,3d\,4p\,5s\,4d\ldots$ (каждая $s$-оболочка вмещает 2 электрона, $p$ --- 6, $d$ --- 10).

\begin{itemize}
\item \textbf{Цирконий}, $Z = 40$, атомная масса 91.22: $[\mathrm{Kr}]\,4d^2\,5s^2$. Четыре валентных электрона, причём
  $d$-оболочка заполнена лишь на две десятых. Такие металлы называют \textbf{переходными}. Их связь --- не только
  «электронный газ», но и заметная направленная составляющая от $d$-электронов. Отсюда высокая температура
  плавления (1855\degC) и склонность к образованию гидридов, оксидов, интерметаллидов.
\item \textbf{Водород}, $Z = 1$: один электрон $1s^1$. В металле он отдаёт свой электрон в общую зону и становится
  экранированным протоном --- самой маленькой частицей, способной сидеть в решётке. Поэтому водород растворяется в
  пустотах между атомами металла и очень быстро по ним прыгает (главы~\ref{ch:lattice} и \ref{ch:diffusion}).
\end{itemize}

\textbf{Электроотрицательность} --- склонность атома тянуть к себе электроны связи. У циркония она низкая (1.33 по
Полингу), у водорода --- 2.20. Поэтому в гидриде электронная плотность сдвинута к водороду: связь Zr--H частично
ионная. Но гидрид остаётся металлом --- проводит ток и даёт металлический блеск на шлифе.

\section{Типы связи}

\begin{center}\small
\begin{tabularx}{\linewidth}{>{\raggedright\arraybackslash}p{2.6cm}>{\raggedright\arraybackslash}X>{\raggedright\arraybackslash}p{3.3cm}}
\toprule
связь & как устроена & примеры \\
\midrule
ионная & электрон переходит от одного атома к другому; ионы притягиваются кулоновски; ненаправленная, но заряды
должны чередоваться & NaCl, ZrO$_2$ (частично) \\
ковалентная & пара электронов общая для двух соседей; связь направлена, число соседей мало & алмаз, Si, SiC \\
металлическая & валентные электроны обобществлены во всём кристалле; положительные ионы в «электронном газе»;
связь ненаправленная & Al, Cu, Zr (с примесью направленной $d$-связи) \\
ван-дер-ваальсова & слабое притяжение наведённых диполей & между цепями полимеров, в графите \\
\bottomrule
\end{tabularx}
\end{center}

Ненаправленность металлической связи объясняет, почему металлы пластичны: атомный слой можно сдвинуть на соседнюю
позицию, не разрывая связей, --- окружение каждого атома после сдвига такое же. В ковалентном кристалле сдвиг рвёт
направленные связи, и кристалл скорее расколется, чем потечёт. Гидриды и оксиды циркония --- промежуточный случай:
у них заметная ионная и ковалентная доля, поэтому они твёрже и хрупче металла.

\section{Энергия пары атомов}

Представим два атома на расстоянии $r$. На больших расстояниях они притягиваются (у металла --- через общие
электроны), на малых отталкиваются: электронные оболочки не могут перекрываться (принцип Паули). Сумма даёт
потенциальную яму (рис.~\ref{fig:a1_potential}а). Удобная модель --- \textbf{потенциал Морзе}
\begin{equation}
U(r) = E_{\text{св}}\Bigl[\bigl(1 - e^{-\alpha(r-r_0)}\bigr)^2 - 1\Bigr].
\end{equation}
Три параметра имеют прямой смысл: $r_0$ --- равновесное расстояние (дно ямы, где сила $F = -\dd U/\dd r$ равна
нулю), $E_{\text{св}}$ --- глубина ямы (энергия связи), $\alpha$ --- обратная ширина ямы.

\fig{a1_potential}{а) Энергия пары атомов (потенциал Морзе) и сила между ними. В минимуме $r_0$ сила равна
нулю; глубина ямы --- энергия связи. б) Уровни энергии колебаний в этой яме и середины размаха колебаний: из-за
несимметричности ямы середина с ростом энергии смещается вправо.}

Из одной кривой $U(r)$ следуют три свойства.
\begin{itemize}
\item \textbf{Температура плавления и теплота испарения} растут с глубиной ямы: чтобы разорвать связи, нужно
  сообщить энергию порядка $E_{\text{св}}$ на атом.
\item \textbf{Модуль упругости} --- это кривизна ямы у дна. Малое удлинение связи $x = r - r_0$ стоит энергии
  $\tfrac12 U''(r_0)x^2$; сила $F = -U''(r_0)\,x$, напряжение --- сила на площадь $r_0^2$, деформация --- $x/r_0$.
  Отсюда
  \begin{equation}
  E \sim \frac{U''(r_0)}{r_0}.
  \end{equation}
  Глубокая и узкая яма --- жёсткий материал.
\item \textbf{Тепловое расширение} --- следствие \emph{несимметричности} ямы (следующий раздел).
\end{itemize}

\begin{idea}
Связь в металле похожа на пружину, но пружину нелинейную: сжимать её много труднее, чем растягивать. Пока
колебания малы, этого не видно, и упругость линейна (закон Гука). Нелинейность проявляется при нагреве (тепловое
расширение) и при очень больших деформациях (предельная прочность связи --- максимум кривой $F(r)$).
\end{idea}

\section{Тепловое расширение}

Атом в яме колеблется. Чем выше температура, тем больше энергия колебаний и размах. В параболической яме
середина размаха не зависит от энергии, и кристалл не расширялся бы. В реальной яме левая стенка круче правой:
при той же энергии атом уходит вправо дальше, чем влево, и средняя длина связи растёт с температурой
(рис.~\ref{fig:a1_potential}б, \ref{fig:a2_thermal}а). Это и есть тепловое расширение. Его коэффициент
\begin{equation}
\alpha_T = \frac{1}{L}\frac{\dd L}{\dd T}
\end{equation}
у металлов порядка $10^{-5}$~К$^{-1}$, и он тем меньше, чем глубже яма: у тугоплавких металлов расширение малое.

\fig{a2_thermal}{а) Средняя длина связи от температуры при классическом (больцмановском) распределении в яме
Морзе и в параболе той же кривизны: в параболе расширения нет. б) Теплоёмкость по Дебаю с $\Theta_D = 290$~К
(цирконий): при низких $T$ закон $T^3$, выше $\Theta_D$ --- предел Дюлонга--Пти $3R$.}

\begin{zrbox}
Цирконий расширяется \emph{анизотропно}: вдоль оси $c$ примерно вдвое сильнее, чем в базисной плоскости (порядка
$10\cdot10^{-6}$ и $5$--$6\cdot10^{-6}$~К$^{-1}$). При охлаждении от температуры отжига соседние зёрна с разной
ориентацией сжимаются по-разному и нагружают друг друга. Так возникают межзёренные напряжения, о которых пойдёт
речь в главе~\ref{ch:plasticity}. Гидрид и матрица тоже расширяются по-разному: несоответствие гидрида
(глава~\ref{ch:elastic}) слегка зависит от температуры.
\end{zrbox}

\section{Колебания решётки и теплоёмкость}

В кристалле атомы колеблются не поодиночке, а согласованно --- волнами. Квант такой волны называют
\textbf{фононом}. Каждый атом имеет три степени свободы колебаний; по классической теореме о равнораспределении на
каждую приходится $k_BT$ (половина --- кинетическая, половина --- потенциальная энергия). Поэтому теплоёмкость моля
атомов
\begin{equation}
C_V = 3R \approx 25~\text{Дж/(моль·К)}
\end{equation}
--- \textbf{закон Дюлонга--Пти}. Для циркония это $25/91.22\approx0.27$~Дж/(г·К).

При низких температурах закон нарушается: высокочастотные колебания «замерзают», потому что их квант
$\hbar\omega$ больше $k_BT$. Модель Дебая описывает это одним параметром --- \textbf{температурой Дебая}
$\Theta_D$ (рис.~\ref{fig:a2_thermal}б). Ниже $\Theta_D$ теплоёмкость падает как $T^3$, выше --- выходит на $3R$. У
циркония $\Theta_D\approx290$~К, поэтому во всём интересном для гидридов диапазоне (комнатная температура и выше)
теплоёмкость практически классическая.

\begin{zrbox}
Теплоёмкость нужна, чтобы читать кривые ДСК (глава~\ref{ch:methods}): прибор меряет тепловой поток, и на
плавном фоне $C_p\,\dd T/\dd t$ выделяются пики растворения и выпадения гидрида. Теплота растворения водорода
--- десятки кДж на моль H, а водорода в образце --- доли атомного процента; поэтому пики невелики на фоне
теплоёмкости (упражнение~\ref{ch:atoms}.2), и фон надо вычитать аккуратно.
\end{zrbox}

\section{Электроны в кристалле: зоны}

В отдельном атоме электрон занимает дискретные уровни. Сблизим $N$ атомов. Волновые функции соседей
перекрываются, и каждый атомный уровень расщепляется на $N$ близких уровней --- образуется \textbf{зона}
(рис.~\ref{fig:a3_bands}а). Ширина зоны растёт с перекрытием, то есть с уменьшением расстояния между атомами.
Между зонами могут оставаться \textbf{запрещённые щели}.

Тот же результат получается и с другой стороны --- от свободного электрона. Электрон с волновым числом $k$
имеет энергию $\hbar^2k^2/2m$ (парабола). В периодической решётке с периодом $a$ волна с $k = \pi/a$ отражается
от атомных плоскостей --- это условие Брэгга (глава~\ref{ch:methods}) для электрона. Стоячие волны с узлами на
атомах и между атомами имеют разную энергию, и на параболе открывается щель
(рис.~\ref{fig:a3_bands}б, модель Кронига--Пенни).

\fig{a3_bands}{а) Сближение атомов: каждый атомный уровень расширяется в зону, ширина растёт с перекрытием
волновых функций (модель сильной связи). б) Электрон в периодическом потенциале (модель Кронига--Пенни): по
сравнению со свободным электроном (пунктир) на $k = n\pi/a$ открываются запрещённые зоны.}

\section{Металлы, полупроводники, диэлектрики}

Электроны заполняют зоны снизу вверх, по два на каждое состояние (принцип Паули). Граница заполнения при $T = 0$
--- \textbf{уровень Ферми} $E_F$. Всё решает, где он оказался (рис.~\ref{fig:a4_fermi}а).
\begin{itemize}
\item Если $E_F$ внутри зоны --- над заполненными состояниями сразу есть свободные, и электрон сколь угодно малым
  полем можно ускорить: это \textbf{металл}.
\item Если зона заполнена целиком, а следующая отделена щелью $E_g$, то при малой щели (около 1~эВ) часть
  электронов перебрасывается теплом --- \textbf{полупроводник}; при большой (несколько эВ) --- \textbf{диэлектрик}.
  ZrO$_2$ --- диэлектрик со щелью около 5~эВ, поэтому оксидная плёнка на оболочке изолирует металл.
\end{itemize}
При $T > 0$ граница заполнения размыта на несколько $k_BT$ (распределение Ферми--Дирака,
рис.~\ref{fig:a4_fermi}б):
\begin{equation}
f(E) = \frac{1}{1 + e^{(E - E_F)/k_BT}}.
\end{equation}
Даже при 1000~К размытие ($k_BT\approx0.09$~эВ) мало по сравнению с шириной зоны (несколько эВ). Поэтому
электроны металла почти не участвуют в теплоёмкости --- она фононная, как и сказано выше.

\fig{a4_fermi}{а) Заполнение зон (закрашено) у металла, полупроводника и диэлектрика. б) Заселённость состояний
по Ферми--Дираку при 0, 300 и 1000~К: размытие порядка $k_BT$.}

\begin{zrbox}
Что из зонной картины нужно для водорода в цирконии.
\begin{itemize}
\item Растворённый водород отдаёт электрон в зону металла; протон экранирован электронами за доли нанометра. В
  решётке он ведёт себя как маленький упругий «шарик», расширяющий окружение: парциальный объём
  $\bar V_{\mathrm H}\approx1.7$~см$^3$/моль (глава~\ref{ch:diffusion}).
\item Гидрид устойчив потому, что связи Zr--H понижают энергию электронов (заполняются связывающие состояния
  ниже $E_F$). Теплота образования гидрида из металла и газа H$_2$ --- порядка $-80$~кДж на моль H, и она примерно
  на 40~кДж/моль больше по модулю, чем теплота растворения водорода в $\alpha$-Zr. Эта разница и есть теплота
  растворения гидрида, задающая наклон линии растворимости (глава~\ref{ch:solutions}).
\item Атомные расчёты (теория функционала плотности) дают из той же электронной структуры упругие модули,
  энергии межфазных границ и барьеры прыжков водорода --- исходные данные для моделей вроде нашей.
\end{itemize}
\end{zrbox}

\section{Итоги главы}

\begin{itemize}
\item Кривая энергии пары атомов задаёт энергию связи (плавление), кривизну у дна (модули упругости) и
  несимметричность (тепловое расширение).
\item Теплоёмкость металла выше температуры Дебая --- $3R$; для циркония это верно уже при комнатной температуре.
\item Электроны в кристалле образуют зоны; металл --- это незаполненная зона. Водород в металле --- экранированный
  протон в пустоте решётки.
\end{itemize}

\begin{exercise}
Для потенциала Морзе найдите жёсткость связи $k = U''(r_0)$ и максимальную силу (предельную прочность связи).
При каком $r$ она достигается? Найдите отношение «теоретической прочности» $F_{\max}/r_0^2$ к модулю упругости
$k/r_0$ при $\alpha r_0 = 1.6$.
\end{exercise}

\begin{exercise}
Сколько теплоты нужно, чтобы нагреть 1~г циркония на 100~К? Сравните с теплотой растворения гидрида: в 1~г
металла со 100~ppm водорода гидрид растворяется с поглощением около 40~кДж на моль H.
\end{exercise}

\begin{exercise}
Чему равна заселённость состояния на 0.1~эВ выше уровня Ферми при 300 и 600~К?
\end{exercise}
